\QHeading{第4题　用简便方法计算下面各题。}
\par\noindent


\CalcBlock[方法一：简便计算]{1}{\begin{aligned}[t]&\phantom{={}}\color{black}125\times(8+4)\\&\color{answerred}=125\times8+125\times4\\&\color{answerred}=1000+500\\&\color{answerred}=1500\end{aligned}}
\par
\MethodBlock{方法二：先求和，再按十位和个位拆分}{把接近整十、整百的因数凑整，再补回或减去多算的部分。}{\begin{aligned}[t]&\phantom{={}}125\times (8+4)\\
&=125\times 12\\
&=125\times (10+2)\\
&=125\times 10+125\times 2\\
&=1250+250\\
&=1500\end{aligned}}

\MethodBlock{方法三：先求和，再拆成4与3相乘}{先算括号里的和，再把12拆成4乘3，用结合律先算125乘4。}{\begin{aligned}[t]&\phantom{={}}125\times (8+4)\\
&=125\times 12\\
&=125\times (4\times 3)\\
&=(125\times 4)\times 3\\
&=500\times 3\\
&=1500\end{aligned}}


\CalcBlock[方法一：简便计算]{2}{\begin{aligned}[t]&\phantom{={}}\color{black}36\times201\\&\color{answerred}=36\times(200+1)\\&\color{answerred}=36\times200+36\times1\\&\color{answerred}=7200+36\\&\color{answerred}=7236\end{aligned}}
\MethodBlock{方法二：按数位拆分计算}{把一个因数拆成整十、整百数与余下部分之和。}{\begin{aligned}[t]&\phantom{={}}36\times 201\\
&=201\times (30+6)\\
&=201\times 30+201\times 6\\
&=6030+1206\\
&=7236\end{aligned}}

\CalcBlock[方法一：简便计算]{3}{\begin{aligned}[t]&\phantom{={}}\color{black}480\div(12\times8)\\&\color{answerred}=480\div12\div8\\&\color{answerred}=40\div8\\&\color{answerred}=5\end{aligned}}
\MethodBlock{方法二：先处理括号或乘除，再拆分除数，连续相除}{}{\begin{aligned}[t]&\phantom{={}}480\div (12\times 8)\\
&=480\div 96\\
&=480\div (2\times 48)\\
&=480\div 2\div 48\\
&=240\div 48\\
&=5\end{aligned}}

\MethodBlock{方法三：先处理括号或乘除，再利用商不变的性质}{被除数和除数同时乘同一个非零数，商不变。}{\begin{aligned}[t]&\phantom{={}}480\div (12\times 8)\\
&=480\div 96\\
&=(480\times 5)\div (96\times 5)\\
&=2400\div 480\\
&=5\end{aligned}}

\CalcBlock[方法一：简便计算]{4}{\begin{aligned}[t]&\phantom{={}}\color{black}780-256-144\\&\color{answerred}=780-(256+144)\\&\color{answerred}=780-400\\&\color{answerred}=380\end{aligned}}
\MethodBlock{方法二：按运算顺序直接计算（对照）}{本法只作结果核对；本题要求简便计算，前面的简便解法仍保留。}{\begin{aligned}[t]&\phantom{={}}780-256-144\\
&=524-144\\
&=380\end{aligned}}

\CalcBlock[方法一：简便计算]{5}{\begin{aligned}[t]&\phantom{={}}\color{black}12.5+3.8-4.6\\&\color{answerred}=12.5-(4.6-3.8)\\&\color{answerred}=12.5-0.8\\&\color{answerred}=11.7\end{aligned}}
\MethodBlock{方法二：按运算顺序直接计算（对照）}{本法只作结果核对；本题要求简便计算，前面的简便解法仍保留。}{\begin{aligned}[t]&\phantom{={}}12.5+3.8-4.6\\
&=16.3-4.6\\
&=11.7\end{aligned}}

\CalcBlock[方法一：简便计算]{6}{\begin{aligned}[t]&\phantom{={}}\color{black}45.6-(12.3+17.7)\\&\color{answerred}=45.6-30\\&\color{answerred}=15.6\end{aligned}}
\MethodBlock{方法二：去括号后计算}{加号后的括号去掉不变号，减号后的括号去掉要改变各项的加减号。}{\begin{aligned}[t]&\phantom{={}}45.6-(12.3+17.7)\\
&=45.6-12.3-17.7\\
&=33.3-17.7\\
&=15.6\end{aligned}}

\par\noindent

\QHeading{第5题　在$\bigcirc$里填上“$>$”“$<$”或“$=$”。}

\PlainItem{$7.3$米 $\textcolor{answerred}{>}$ $7$米$3$厘米}
\par
\PlainItem{$6$平方米 $\textcolor{answerred}{>}$ $590$平方分米}
\PlainItem{$3.6$角 $\textcolor{answerred}{=}$ $3.60$角}
\PlainItem{$2$吨$10$千克 $\textcolor{answerred}{<}$ $2.1$吨}
\PlainItem{$9$千米$200$米 $\textcolor{answerred}{=}$ $9.2$千米}
\PlainItem{$3$平方米$2$平方分米 $\textcolor{answerred}{<}$ $3.2$平方米}
\par\noindent

\Answer{
\QHeading{计算过程：第5题：统一单位后比较大小}
\begin{enumerate}[itemsep=4pt]
\item 7米3厘米化成米：
\Formula{
\begin{aligned}[t]&\phantom{={}}7+3\div100\\&=7+0.03\\&=7.03\text{（米）}\end{aligned}
}
$7.3>7.03$，填$>$。
\item 6平方米化成平方分米：
\Formula{
\begin{aligned}[t]&\phantom{={}}6\times100\\&=600\text{（平方分米）}\end{aligned}
}
$600>590$，填$>$。
\item 3.6角与3.60角：
小数末尾添0，大小不变，填$=$。
\item 2吨10千克化成吨：
\Formula{
\begin{aligned}[t]&\phantom{={}}2+10\div1000\\&=2+0.01\\&=2.01\text{（吨）}\end{aligned}
}
$2.01<2.1$，填$<$。
\item 9千米200米化成千米：
\Formula{
\begin{aligned}[t]&\phantom{={}}9+200\div1000\\&=9+0.2\\&=9.2\text{（千米）}\end{aligned}
}
$9.2=9.2$，填$=$。
\item 3平方米2平方分米化成平方米：
\Formula{
\begin{aligned}[t]&\phantom{={}}3+2\div100\\&=3+0.02\\&=3.02\text{（平方米）}\end{aligned}
}
$3.02<3.2$，填$<$。
\end{enumerate}

}

\QHeading{第6题　分别画出下面图形底边上的高。}

\par\noindent\makebox[\linewidth][c]{
\begin{tikzpicture}[x=1bp,y=-1bp,color=black,draw=black,text=black]
\begin{scope}\clip (56.0,481.0) rectangle (522.0,578.0);
\hdhead{98}{综合测试}
\hdsection{4}{147}{用简便方法计算下面各题。}
\q{106}{191}{257}{12}{$125\times(8+4)$}
\q{354}{191}{257}{12}{$36\times201$}
\q{598}{191}{257}{12}{$480\div(12\times8)$}
\q{106}{335}{257}{12}{$780-256-144$}
\q{354}{335}{257}{12}{$12.5+3.8-4.6$}
\q{598}{335}{257}{12}{$45.6-(12.3+17.7)$}
\hdsection{5}{492}{在$\bigcirc$里填上“$>$”“$<$”或“$=$”。}
\q{101}{536}{80}{12}{$7.3$米}
\q{211}{536}{200}{12}{$7$米$3$厘米}
\draw[line width=.6pt] (123.333,364.000) circle (9.333bp);
\q{101}{581}{80}{12}{$3.6$角}
\q{211}{581}{200}{12}{$3.60$角}
\draw[line width=.6pt] (123.333,394.000) circle (9.333bp);
\q{101}{626}{155}{12}{$9$千米$200$米}
\q{286}{626}{200}{12}{$9.2$千米}
\draw[line width=.6pt] (173.333,424.000) circle (9.333bp);
\q{440}{536}{104}{12}{$6$平方米}
\q{574}{536}{220}{12}{$590$平方分米}
\draw[line width=.6pt] (365.333,364.000) circle (9.333bp);
\q{440}{581}{139}{12}{$2$吨$10$千克}
\q{609}{581}{220}{12}{$2.1$吨}
\draw[line width=.6pt] (388.667,394.000) circle (9.333bp);
\q{440}{626}{222}{12}{$3$平方米$2$平方分米}
\q{692}{626}{220}{12}{$3.2$平方米}
\draw[line width=.6pt] (444.000,424.000) circle (9.333bp);
\hdsection{6}{687}{分别画出下面图形底边上的高。}
\draw[line width=.7pt] (95.333,487.333)--(116.000,536.667)--(162.000,536.667)--cycle;
\draw[line width=.7pt] (211.333,487.333)--(211.333,536.667)--(259.333,536.667)--cycle;
\draw[line width=.7pt] (298.000,487.333)--(359.333,487.333)--(386.667,536.667)--(324.000,536.667)--cycle;
\draw[line width=.7pt] (433.333,487.333)--(506.667,487.333)--(516.000,537.333)--(424.000,537.333)--cycle;
\draw[line width=.5pt] (211.333,530.000)--(218.000,530.000)--(218.000,536.667);
\node[inner sep=0pt,align=center,font=\normalsize,] at (140.667,549.333) {底};
\node[inner sep=0pt,align=center,font=\normalsize,] at (244.667,504.000) {底};
\node[inner sep=0pt,align=center,font=\normalsize,] at (354.667,549.333) {底};
\node[inner sep=0pt,align=center,font=\normalsize,] at (467.333,549.333) {底};
\hdsection{7}{870}{按要求画一画。}
\q{97}{913}{354}{12}{① 将图形绕$O$点逆时针旋转$90^\circ$，}
\q{129}{954}{330}{12}{再向左平移5格。}
\q{471}{913}{354}{12}{② 将三角形绕$O$点顺时针旋}
\q{505}{954}{315}{12}{转$90^\circ$。}
\draw[hdgrid] (110.000,658.000)--(110.000,769.333);
\draw[hdgrid] (123.917,658.000)--(123.917,769.333);
\draw[hdgrid] (137.833,658.000)--(137.833,769.333);
\draw[hdgrid] (151.750,658.000)--(151.750,769.333);
\draw[hdgrid] (165.667,658.000)--(165.667,769.333);
\draw[hdgrid] (179.583,658.000)--(179.583,769.333);
\draw[hdgrid] (193.500,658.000)--(193.500,769.333);
\draw[hdgrid] (207.417,658.000)--(207.417,769.333);
\draw[hdgrid] (221.333,658.000)--(221.333,769.333);
\draw[hdgrid] (235.250,658.000)--(235.250,769.333);
\draw[hdgrid] (110.000,658.000)--(235.250,658.000);
\draw[hdgrid] (110.000,671.917)--(235.250,671.917);
\draw[hdgrid] (110.000,685.833)--(235.250,685.833);
\draw[hdgrid] (110.000,699.750)--(235.250,699.750);
\draw[hdgrid] (110.000,713.667)--(235.250,713.667);
\draw[hdgrid] (110.000,727.583)--(235.250,727.583);
\draw[hdgrid] (110.000,741.500)--(235.250,741.500);
\draw[hdgrid] (110.000,755.417)--(235.250,755.417);
\draw[hdgrid] (110.000,769.333)--(235.250,769.333);
\draw[hdgrid] (356.667,659.333)--(356.667,770.267);
\draw[hdgrid] (370.533,659.333)--(370.533,770.267);
\draw[hdgrid] (384.400,659.333)--(384.400,770.267);
\draw[hdgrid] (398.267,659.333)--(398.267,770.267);
\draw[hdgrid] (412.133,659.333)--(412.133,770.267);
\draw[hdgrid] (426.000,659.333)--(426.000,770.267);
\draw[hdgrid] (439.867,659.333)--(439.867,770.267);
\draw[hdgrid] (453.733,659.333)--(453.733,770.267);
\draw[hdgrid] (467.600,659.333)--(467.600,770.267);
\draw[hdgrid] (481.467,659.333)--(481.467,770.267);
\draw[hdgrid] (495.333,659.333)--(495.333,770.267);
\draw[hdgrid] (509.200,659.333)--(509.200,770.267);
\draw[hdgrid] (523.067,659.333)--(523.067,770.267);
\draw[hdgrid] (356.667,659.333)--(523.067,659.333);
\draw[hdgrid] (356.667,673.200)--(523.067,673.200);
\draw[hdgrid] (356.667,687.067)--(523.067,687.067);
\draw[hdgrid] (356.667,700.933)--(523.067,700.933);
\draw[hdgrid] (356.667,714.800)--(523.067,714.800);
\draw[hdgrid] (356.667,728.667)--(523.067,728.667);
\draw[hdgrid] (356.667,742.533)--(523.067,742.533);
\draw[hdgrid] (356.667,756.400)--(523.067,756.400);
\draw[hdgrid] (356.667,770.267)--(523.067,770.267);
\draw[line width=.7pt] (193.333,700.000)--(193.333,742.000)--(220.667,728.000);
\fill (193.333,700.000) circle (1.1bp);
\node[inner sep=0pt,align=center,font=\normalsize,] at (186.667,694.000) {$O$};
\draw[line width=.7pt] (398.667,742.667)--(398.667,659.333)--(426.000,728.667)--cycle;
\fill (398.667,742.667) circle (1.1bp);
\node[inner sep=0pt,align=center,font=\normalsize,] at (392.000,748.667) {$O$};

\begin{scope}[color=answerred,draw=answerred,text=answerred]\ans{71.333}{150.667}{156.667}{10}{\linespread{1}\selectfont\setlength{\jot}{1pt}$\begin{aligned}&=125\times8+125\times4\\&=1000+500\\&=1500\end{aligned}$}
\ans{234.667}{150.667}{156.667}{10}{\linespread{1}\selectfont\setlength{\jot}{1pt}$\begin{aligned}&=36\times(200+1)\\&=36\times200+36\times1\\&=7200+36\\&=7236\end{aligned}$}
\ans{398.667}{150.667}{156.667}{10}{\linespread{1}\selectfont\setlength{\jot}{1pt}$\begin{aligned}&=480\div12\div8\\&=40\div8\\&=5\end{aligned}$}
\ans{71.333}{247.333}{156.667}{10}{\linespread{1}\selectfont\setlength{\jot}{1pt}$\begin{aligned}&=780-(256+144)\\&=780-400\\&=380\end{aligned}$}
\ans{234.667}{247.333}{156.667}{10}{\linespread{1}\selectfont\setlength{\jot}{1pt}$\begin{aligned}&=12.5+3.8-(3.8+0.8)\\&=12.5+3.8-3.8-0.8\\&=12.5+(3.8-3.8)-0.8\\&=12.5+0-0.8\\&=12.5-0.8\\&=11.7\end{aligned}$}
\ans{398.667}{247.333}{156.667}{10}{\linespread{1}\selectfont\setlength{\jot}{1pt}$\begin{aligned}&=45.6-30\\&=15.6\end{aligned}$}
\ans{118.000}{359.333}{30.000}{13}{$>$}
\ans{360.667}{359.333}{30.000}{13}{$>$}
\ans{118.000}{389.333}{30.000}{13}{$=$}
\ans{383.333}{389.333}{30.000}{13}{$<$}
\ans{165.333}{420.667}{30.000}{13}{$=$}
\ans{441.333}{420.667}{30.000}{13}{$<$}
\begin{scope}[scale=0.6666667]\draw[answerred,dashed,line width=.9pt] (143,731)--(143,805);\draw[answerred,dashed](143,805)--(175,805);\draw (143,796)--(152,796)--(152,805);\node[left]at(143,768){高};\draw[answerred,dashed,line width=.9pt](316.9995,805.0005)--(353.9861,769.0140);\draw[answerred](348.4073,763.2802)--(342.6735,768.8590)--(348.2523,774.5928);\node[left,inner sep=2pt] at(335.4928,787.0073){高};\draw[answerred,dashed,line width=.9pt](539,731)--(539,805);\draw(539,796)--(548,796)--(548,805);\node[left]at(539,767){高};\draw[answerred,dashed,line width=.9pt](650,731)--(650,806);\draw(650,797)--(659,797)--(659,806);\node[right]at(650,766){高};\end{scope}
\begin{scope}[scale=0.6666667]\draw[answerred,dashed,line width=.75pt](290,1050)--(352,1050)--(331,1009);\draw[answerred,line width=1.1pt](186,1050)--(248,1050)--(228,1009);\fill[answerred](186,1050)circle(1.6pt);\node[below left]at(186,1050){$O\!\prime$};\draw[answerred,line width=1.1pt](598,1114)--(722,1114)--(619,1155)--cycle;\end{scope}
\ans{150.000}{787.000}{400.000}{9}{完整计算过程见紧接本页的计算续页。}
\end{scope}
\end{scope}\pgfresetboundingbox\path[use as bounding box] (56.0,481.0) rectangle (522.0,578.0);
\end{tikzpicture}}\par



\par







\par\noindent

\QHeading{第7题　按要求画一画。}

\par\noindent\makebox[\linewidth][c]{
\begin{tikzpicture}[x=1bp,y=-1bp,color=black,draw=black,text=black]
\begin{scope}\clip (59.0,603.0) rectangle (528.0,776.0);
\hdhead{98}{综合测试}
\hdsection{4}{147}{用简便方法计算下面各题。}
\q{106}{191}{257}{12}{$125\times(8+4)$}
\q{354}{191}{257}{12}{$36\times201$}
\q{598}{191}{257}{12}{$480\div(12\times8)$}
\q{106}{335}{257}{12}{$780-256-144$}
\q{354}{335}{257}{12}{$12.5+3.8-4.6$}
\q{598}{335}{257}{12}{$45.6-(12.3+17.7)$}
\hdsection{5}{492}{在$\bigcirc$里填上“$>$”“$<$”或“$=$”。}
\q{101}{536}{80}{12}{$7.3$米}
\q{211}{536}{200}{12}{$7$米$3$厘米}
\draw[line width=.6pt] (123.333,364.000) circle (9.333bp);
\q{101}{581}{80}{12}{$3.6$角}
\q{211}{581}{200}{12}{$3.60$角}
\draw[line width=.6pt] (123.333,394.000) circle (9.333bp);
\q{101}{626}{155}{12}{$9$千米$200$米}
\q{286}{626}{200}{12}{$9.2$千米}
\draw[line width=.6pt] (173.333,424.000) circle (9.333bp);
\q{440}{536}{104}{12}{$6$平方米}
\q{574}{536}{220}{12}{$590$平方分米}
\draw[line width=.6pt] (365.333,364.000) circle (9.333bp);
\q{440}{581}{139}{12}{$2$吨$10$千克}
\q{609}{581}{220}{12}{$2.1$吨}
\draw[line width=.6pt] (388.667,394.000) circle (9.333bp);
\q{440}{626}{222}{12}{$3$平方米$2$平方分米}
\q{692}{626}{220}{12}{$3.2$平方米}
\draw[line width=.6pt] (444.000,424.000) circle (9.333bp);
\hdsection{6}{687}{分别画出下面图形底边上的高。}
\draw[line width=.7pt] (95.333,487.333)--(116.000,536.667)--(162.000,536.667)--cycle;
\draw[line width=.7pt] (211.333,487.333)--(211.333,536.667)--(259.333,536.667)--cycle;
\draw[line width=.7pt] (298.000,487.333)--(359.333,487.333)--(386.667,536.667)--(324.000,536.667)--cycle;
\draw[line width=.7pt] (433.333,487.333)--(506.667,487.333)--(516.000,537.333)--(424.000,537.333)--cycle;
\draw[line width=.5pt] (211.333,530.000)--(218.000,530.000)--(218.000,536.667);
\node[inner sep=0pt,align=center,font=\normalsize,] at (140.667,549.333) {底};
\node[inner sep=0pt,align=center,font=\normalsize,] at (244.667,504.000) {底};
\node[inner sep=0pt,align=center,font=\normalsize,] at (354.667,549.333) {底};
\node[inner sep=0pt,align=center,font=\normalsize,] at (467.333,549.333) {底};
\hdsection{7}{870}{按要求画一画。}
\q{97}{913}{354}{12}{① 将图形绕$O$点逆时针旋转$90^\circ$，}
\q{129}{954}{330}{12}{再向左平移5格。}
\q{471}{913}{354}{12}{② 将三角形绕$O$点顺时针旋}
\q{505}{954}{315}{12}{转$90^\circ$。}
\draw[hdgrid] (110.000,658.000)--(110.000,769.333);
\draw[hdgrid] (123.917,658.000)--(123.917,769.333);
\draw[hdgrid] (137.833,658.000)--(137.833,769.333);
\draw[hdgrid] (151.750,658.000)--(151.750,769.333);
\draw[hdgrid] (165.667,658.000)--(165.667,769.333);
\draw[hdgrid] (179.583,658.000)--(179.583,769.333);
\draw[hdgrid] (193.500,658.000)--(193.500,769.333);
\draw[hdgrid] (207.417,658.000)--(207.417,769.333);
\draw[hdgrid] (221.333,658.000)--(221.333,769.333);
\draw[hdgrid] (235.250,658.000)--(235.250,769.333);
\draw[hdgrid] (110.000,658.000)--(235.250,658.000);
\draw[hdgrid] (110.000,671.917)--(235.250,671.917);
\draw[hdgrid] (110.000,685.833)--(235.250,685.833);
\draw[hdgrid] (110.000,699.750)--(235.250,699.750);
\draw[hdgrid] (110.000,713.667)--(235.250,713.667);
\draw[hdgrid] (110.000,727.583)--(235.250,727.583);
\draw[hdgrid] (110.000,741.500)--(235.250,741.500);
\draw[hdgrid] (110.000,755.417)--(235.250,755.417);
\draw[hdgrid] (110.000,769.333)--(235.250,769.333);
\draw[hdgrid] (356.667,659.333)--(356.667,770.267);
\draw[hdgrid] (370.533,659.333)--(370.533,770.267);
\draw[hdgrid] (384.400,659.333)--(384.400,770.267);
\draw[hdgrid] (398.267,659.333)--(398.267,770.267);
\draw[hdgrid] (412.133,659.333)--(412.133,770.267);
\draw[hdgrid] (426.000,659.333)--(426.000,770.267);
\draw[hdgrid] (439.867,659.333)--(439.867,770.267);
\draw[hdgrid] (453.733,659.333)--(453.733,770.267);
\draw[hdgrid] (467.600,659.333)--(467.600,770.267);
\draw[hdgrid] (481.467,659.333)--(481.467,770.267);
\draw[hdgrid] (495.333,659.333)--(495.333,770.267);
\draw[hdgrid] (509.200,659.333)--(509.200,770.267);
\draw[hdgrid] (523.067,659.333)--(523.067,770.267);
\draw[hdgrid] (356.667,659.333)--(523.067,659.333);
\draw[hdgrid] (356.667,673.200)--(523.067,673.200);
\draw[hdgrid] (356.667,687.067)--(523.067,687.067);
\draw[hdgrid] (356.667,700.933)--(523.067,700.933);
\draw[hdgrid] (356.667,714.800)--(523.067,714.800);
\draw[hdgrid] (356.667,728.667)--(523.067,728.667);
\draw[hdgrid] (356.667,742.533)--(523.067,742.533);
\draw[hdgrid] (356.667,756.400)--(523.067,756.400);
\draw[hdgrid] (356.667,770.267)--(523.067,770.267);
\draw[line width=.7pt] (193.333,700.000)--(193.333,742.000)--(220.667,728.000);
\fill (193.333,700.000) circle (1.1bp);
\node[inner sep=0pt,align=center,font=\normalsize,] at (186.667,694.000) {$O$};
\draw[line width=.7pt] (398.667,742.667)--(398.667,659.333)--(426.000,728.667)--cycle;
\fill (398.667,742.667) circle (1.1bp);
\node[inner sep=0pt,align=center,font=\normalsize,] at (392.000,748.667) {$O$};

\begin{scope}[color=answerred,draw=answerred,text=answerred]\ans{71.333}{150.667}{156.667}{10}{\linespread{1}\selectfont\setlength{\jot}{1pt}$\begin{aligned}&=125\times8+125\times4\\&=1000+500\\&=1500\end{aligned}$}
\ans{234.667}{150.667}{156.667}{10}{\linespread{1}\selectfont\setlength{\jot}{1pt}$\begin{aligned}&=36\times(200+1)\\&=36\times200+36\times1\\&=7200+36\\&=7236\end{aligned}$}
\ans{398.667}{150.667}{156.667}{10}{\linespread{1}\selectfont\setlength{\jot}{1pt}$\begin{aligned}&=480\div12\div8\\&=40\div8\\&=5\end{aligned}$}
\ans{71.333}{247.333}{156.667}{10}{\linespread{1}\selectfont\setlength{\jot}{1pt}$\begin{aligned}&=780-(256+144)\\&=780-400\\&=380\end{aligned}$}
\ans{234.667}{247.333}{156.667}{10}{\linespread{1}\selectfont\setlength{\jot}{1pt}$\begin{aligned}&=12.5+3.8-(3.8+0.8)\\&=12.5+3.8-3.8-0.8\\&=12.5+(3.8-3.8)-0.8\\&=12.5+0-0.8\\&=12.5-0.8\\&=11.7\end{aligned}$}
\ans{398.667}{247.333}{156.667}{10}{\linespread{1}\selectfont\setlength{\jot}{1pt}$\begin{aligned}&=45.6-30\\&=15.6\end{aligned}$}
\ans{118.000}{359.333}{30.000}{13}{$>$}
\ans{360.667}{359.333}{30.000}{13}{$>$}
\ans{118.000}{389.333}{30.000}{13}{$=$}
\ans{383.333}{389.333}{30.000}{13}{$<$}
\ans{165.333}{420.667}{30.000}{13}{$=$}
\ans{441.333}{420.667}{30.000}{13}{$<$}
\begin{scope}[scale=0.6666667]\draw[answerred,dashed,line width=.9pt] (143,731)--(143,805);\draw[answerred,dashed](143,805)--(175,805);\draw (143,796)--(152,796)--(152,805);\node[left]at(143,768){高};\draw[answerred,dashed,line width=.9pt](316.9995,805.0005)--(353.9861,769.0140);\draw[answerred](348.4073,763.2802)--(342.6735,768.8590)--(348.2523,774.5928);\node[left,inner sep=2pt] at(335.4928,787.0073){高};\draw[answerred,dashed,line width=.9pt](539,731)--(539,805);\draw(539,796)--(548,796)--(548,805);\node[left]at(539,767){高};\draw[answerred,dashed,line width=.9pt](650,731)--(650,806);\draw(650,797)--(659,797)--(659,806);\node[right]at(650,766){高};\end{scope}
\begin{scope}[scale=0.6666667]\draw[answerred,dashed,line width=.75pt](290,1050)--(352,1050)--(331,1009);\draw[answerred,line width=1.1pt](186,1050)--(248,1050)--(228,1009);\fill[answerred](186,1050)circle(1.6pt);\node[below left]at(186,1050){$O\!\prime$};\draw[answerred,line width=1.1pt](598,1114)--(722,1114)--(619,1155)--cycle;\end{scope}

\end{scope}
\end{scope}\pgfresetboundingbox\path[use as bounding box] (59.0,603.0) rectangle (528.0,776.0);
\end{tikzpicture}}\par



\par







\QHeading{第6、7题　不同作法与不同答案的区别}

\Answer{三角形给定底边后，对应的高唯一，可用三角板或量角器作出；平行四边形与梯形的高可选不同位置。给定中心、角度和方向的旋转结果唯一，逐点数格法和量角器法应得到相同结果。不要为了增加“做法”改变题目规定的底边、中心、角度或方向。}

\Answer{作高复核：高用虚线，旁边写“高”，垂足处画直角符号；以题目指定的底为准，必要时用虚线延长底边。}
